%% ==========================================================================
%%  Unit Subsidies for Negative Dichotomous Valuations
%%
%%  The main theorem, for an arbitrary number of agents.  Promoted from
%%  working/approach_13.tex.  Report form: no motivation, no status
%%  paragraphs, no computational verification, no discussion of routes not
%%  taken.  Every claim is proved.
%%
%%  ATTRIBUTION.  Theorem~\ref{thm:m-twyz} and the algorithm described in
%%  \Cref{sec:m-terminal} are due to Tao, Wu, Yu and Zhou~\cite{TWYZ23}; they
%%  are restated here, not claimed.  Everything from
%%  \Cref{sec:m-completion} onwards is new.
%%
%%  REFERENCING: theorem-like environments share one counter in this
%%  preamble, so \Cref renders all of them as "Theorem".  Use explicit
%%  "Lemma~\ref{...}", "Theorem~\ref{...}", "Definition~\ref{...}".
%% ==========================================================================
\section{Unit Subsidies for Negative Dichotomous Valuations}
\label{sec:main}

This section proves the main theorem: for every number of agents $n$ and every
negative dichotomous instance there is an allocation admitting an envy-free
solution with a subsidy of $0$ or $1$ per agent, at most $n-1$ in total,
computable in polynomial time. Beyond Definition~\ref{def:dichcost} nothing is
assumed of the cost functions --- not additivity, not submodularity, not any
relation between different agents' cost functions --- and the number of chores
is unrestricted.

The argument takes as input an envy-free \emph{partial} allocation produced by
the algorithm of Tao, Wu, Yu and Zhou~\cite{TWYZ23}, and completes it. That
algorithm, its Theorem~5.1, and the three update rules recorded in
\Cref{sec:m-terminal} are due to~\cite{TWYZ23} and are restated here for
reference; the completion of \Cref{sec:m-completion} onwards, and every lemma
and theorem after Lemma~\ref{lem:m-terminal}, are ours.

We work in cost form throughout, so an envy-free solution
(Definition~\ref{def:efsolution}) is the requirement
\begin{equation}
\label{eq:m-ef}
  \cost_i(\bundle i) - \subsidy_i \;\le\; \cost_i(\bundle j) - \subsidy_j
  \qquad \text{for all } i,j \in \agents ,
\end{equation}
and a partial allocation is envy-free (Definition~\ref{def:ef}) when
\begin{equation}
\label{eq:m-pef}
  \cost_i(\bundle i) \;\le\; \cost_i(\bundle j)
  \qquad \text{for all } i,j \in \agents .
\end{equation}
For $S \subseteq \items$ and $e \in \items \setminus S$ write
$\cost_i(e \mid S) := \cost_i(S \cup \set e) - \cost_i(S) \in \set{0,1}$ for a
marginal, and for a partial allocation $\alloc$ write
$R(\alloc) := \items \setminus \bigcup_{i \in \agents} \bundle i$ for its
\emph{residue}.

\begin{observation}
\label{obs:m-gap}
For every $i \in \agents$ and all $S,T \subseteq \items$,
\[
  \cost_i(S) < \cost_i(T) \quad\Longrightarrow\quad \cost_i(S) \le \cost_i(T) - 1 .
\]
Moreover $\cost_i(S) \le \cost_i(T)$ whenever $S \subseteq T$.
\end{observation}

\begin{proof}
By Observation~\ref{obs:integrality} every $\cost_i$ is integer valued, so a
strict inequality between two of its values is a gap of at least $1$. For
monotonicity, enumerate $T \setminus S = \set{f_1,\dots,f_q}$ and telescope:
$\cost_i(T) - \cost_i(S) = \sum_{u=1}^{q} \cost_i\bigl(f_u \mid S \cup
\set{f_1,\dots,f_{u-1}}\bigr)$, a sum of terms in $\set{0,1}$ by
Definition~\ref{def:dichcost}, hence non-negative.
\end{proof}

% --------------------------------------------------------------------------
\subsection{The terminal state of the partial-allocation algorithm}
\label{sec:m-terminal}

\begin{theorem}[\cite{TWYZ23}, Theorem~5.1]
\label{thm:m-twyz}
For every $n$ and all dichotomous cost functions $\cost_1,\dots,\cost_n$ there
is a polynomial-time algorithm returning a partial allocation $X$ that is
envy-free in the sense of \eqref{eq:m-pef} and satisfies
$\abs{R(X)} \le n-1$.
\end{theorem}

The proof below uses not only the statement of Theorem~\ref{thm:m-twyz} but the
state in which the algorithm of~\cite{TWYZ23} halts, so we record that
algorithm. It maintains a partial allocation $X = (X_1,\dots,X_n)$, envy-free
throughout, together with the \emph{equality graph}
\begin{equation}
\label{eq:m-graph}
  G = (\agents, E),
  \qquad
  (i,j) \in E \iff i \ne j \ \text{ and } \ \cost_i(X_i) = \cost_i(X_j),
\end{equation}
and repeats the following three rules in the stated order, halting only when
none applies.
\begin{enumerate}
  \item[(R1)] If some $e \in R(X)$ and some $i \in \agents$ satisfy
        $\cost_i(e \mid X_i) = 0$, assign $e$ to $i$.
  \item[(R2)] If some $e \in R(X)$ and some arc $(i,j) \in E$ lying on a
        directed cycle of $G$ satisfy $\cost_i(e \mid X_j) = 0$, rotate the
        bundles along that cycle and assign $e$ to $i$.
  \item[(R3)] Otherwise select a strongly connected component $S$ of $G$ with
        no arc of $E$ leaving $S$. If $\abs{R(X)} \ge \abs S$, assign one
        residual chore to each agent of $S$; otherwise halt.
\end{enumerate}
The loop runs while $R(X) \ne \emptyset$, so the only exit with
$R(X) \ne \emptyset$ is the final clause of (R3).

For the remainder of this section fix the following data: let $X$ be the
partial allocation at which the algorithm halts, let $R := R(X)$ and
$r := \abs R$, and, when $r \ge 1$, let $S$ be the component selected by (R3)
at the halting step, with $G$ and $E$ as in \eqref{eq:m-graph} evaluated at
that state.

\begin{lemma}
\label{lem:m-terminal}
Suppose $r \ge 1$. Then
\begin{enumerate}
  \item[(a)] $r < \abs S$;
  \item[(b)] $\cost_i(e \mid X_i) = 1$ for every $i \in \agents$ and every
        $e \in R$;
  \item[(c)] $\cost_i(e \mid X_j) = 1$ for every arc $(i,j) \in E$ with
        $i,j \in S$ and every $e \in R$;
  \item[(d)] $\cost_i(X_i) < \cost_i(X_j)$ for every $i \in S$ and every
        $j \in \agents \setminus S$.
\end{enumerate}
\end{lemma}

\begin{proof}
Since $r \ge 1$, the algorithm halted at the final clause of (R3); hence
neither (R1) nor (R2) applied at the halting state, and $\abs R < \abs S$,
which is (a).

(b) If $\cost_i(e \mid X_i) = 0$ for some $i \in \agents$ and some $e \in R$,
rule (R1) would apply. Marginals lie in $\set{0,1}$ by
Definition~\ref{def:dichcost}, so the value is $1$.

(c) Let $(i,j) \in E$ with $i,j \in S$. Since $S$ is strongly connected there
is a directed path in $G$ from $j$ to $i$; taking one with no repeated vertex
and appending the arc $(i,j)$ yields a directed cycle of $G$ through $(i,j)$.
If $\cost_i(e \mid X_j) = 0$ for some $e \in R$, rule (R2) would apply to that
arc and that cycle. Hence the marginal is not $0$, so it is $1$.

(d) Let $i \in S$ and $j \in \agents \setminus S$. No arc of $E$ leaves $S$, so
$(i,j) \notin E$, i.e.\ $\cost_i(X_i) \ne \cost_i(X_j)$ by \eqref{eq:m-graph}.
Envy-freeness \eqref{eq:m-pef} gives $\cost_i(X_i) \le \cost_i(X_j)$, so the
inequality is strict.
\end{proof}

% --------------------------------------------------------------------------
\subsection{The completion}
\label{sec:m-completion}

Assume $r \ge 1$ and write $R = \set{e_1,\dots,e_r}$. By
Lemma~\ref{lem:m-terminal}(a), $r < \abs S$, so we may fix $r$ distinct agents
\[
  T = \set{t_1,\dots,t_r} \subseteq S ,
\]
chosen arbitrarily. Define $\alloc = (A_1,\dots,A_n)$ by
\begin{equation}
\label{eq:m-alloc}
  A_i :=
  \begin{cases}
    X_{t_k} \cup \set{e_k}, & i = t_k \ \ (1 \le k \le r),\\[2pt]
    X_i, & i \in \agents \setminus T .
  \end{cases}
\end{equation}
The $t_k$ are distinct and each chore of $R$ is used exactly once, so the
bundles of $\alloc$ are pairwise disjoint and cover $\items$: $\alloc$ is a
complete allocation. Members of $T$ are called \emph{recipients}.

\begin{definition}[Subsidy set]
\label{def:m-P}
Let $P \subseteq \agents$ be the smallest set satisfying
\begin{enumerate}
  \item[(P1)] $T \subseteq P$; and
  \item[(P2)] if $i \in \agents \setminus S$ and $(i,j) \in E$ for some
        $j \in P$, then $i \in P$.
\end{enumerate}
Equivalently, $P$ is obtained by initialising $P := T$ and repeatedly adjoining
any $i \in \agents \setminus S$ having an arc of $E$ into the current set,
until no further agent can be adjoined; the process terminates because
$\agents$ is finite and $P$ only grows. Define
\begin{equation}
\label{eq:m-subsidy}
  \subsidy_i := \begin{cases} 1, & i \in P,\\ 0, & i \in \agents \setminus P.\end{cases}
\end{equation}
\end{definition}

\begin{lemma}
\label{lem:m-Pstructure}
$P \setminus T \subseteq \agents \setminus S$: every agent of $P$ that is not a
recipient lies outside $S$.
\end{lemma}

\begin{proof}
Let $P_0 := T$ and let $P_{u+1}$ be $P_u$ together with every
$i \in \agents \setminus S$ having an arc of $E$ into $P_u$, so that
$P = \bigcup_{u \ge 0} P_u$. We show $P_u \setminus T \subseteq \agents
\setminus S$ by induction on $u$. For $u = 0$ the set $P_0 \setminus T$ is
empty. For the step, every agent adjoined in passing from $P_u$ to $P_{u+1}$
lies in $\agents \setminus S$ by construction.
\end{proof}

\begin{lemma}
\label{lem:m-budget}
$\subsidy \in \set{0,1}^n$ and $\sum_{i \in \agents} \subsidy_i = \abs P \le n-1$.
\end{lemma}

\begin{proof}
That $\subsidy \in \set{0,1}^n$ with $\sum_i \subsidy_i = \abs P$ is immediate
from \eqref{eq:m-subsidy}. By Lemma~\ref{lem:m-Pstructure},
$P \subseteq T \cup (\agents \setminus S)$, and this union is disjoint because
$T \subseteq S$; hence
$\abs P \le \abs T + \abs{\agents \setminus S} = r + n - \abs S$. By
Lemma~\ref{lem:m-terminal}(a), $r \le \abs S - 1$, so $\abs P \le n-1$.
\end{proof}

\begin{lemma}
\label{lem:m-expensive}
Let $i \in S$, let $t \in T$ and let $e \in R$. Then
\[
  \cost_i\bigl(X_t \cup \set e\bigr) \;\ge\; \cost_i(X_i) + 1 .
\]
\end{lemma}

\begin{proof}
Note $t \in T \subseteq S$. Envy-freeness \eqref{eq:m-pef} gives
$\cost_i(X_i) \le \cost_i(X_t)$.

If $\cost_i(X_i) < \cost_i(X_t)$, then $\cost_i(X_t) \ge \cost_i(X_i) + 1$ by
Observation~\ref{obs:m-gap}, and $\cost_i(X_t \cup \set e) \ge \cost_i(X_t)$ by
monotonicity, again Observation~\ref{obs:m-gap}.

Otherwise $\cost_i(X_i) = \cost_i(X_t)$. If $i = t$ then
$\cost_i(X_t \cup \set e) = \cost_i(X_i) + \cost_i(e \mid X_i) = \cost_i(X_i)+1$
by Lemma~\ref{lem:m-terminal}(b). If $i \ne t$ then $(i,t) \in E$ by
\eqref{eq:m-graph}, with $i,t \in S$, so $\cost_i(e \mid X_t) = 1$ by
Lemma~\ref{lem:m-terminal}(c) and therefore
$\cost_i(X_t \cup \set e) = \cost_i(X_t) + 1 = \cost_i(X_i) + 1$.
\end{proof}

% --------------------------------------------------------------------------
\subsection{Envy-freeness of the completion}
\label{sec:m-ef}

\begin{theorem}
\label{thm:m-completion}
Let $X$, $R$, $S$ be as in \Cref{sec:m-terminal} with $r \ge 1$, let $\alloc$
be as in \eqref{eq:m-alloc} and let $\subsidy$ be as in \eqref{eq:m-subsidy}.
Then $(\alloc,\subsidy)$ is an envy-free solution.
\end{theorem}

\begin{proof}
Fix $i,j \in \agents$; we verify \eqref{eq:m-ef}, that is
$\cost_i(A_i) - \subsidy_i \le \cost_i(A_j) - \subsidy_j$. Throughout,
$X_j \subseteq A_j$ for every $j$ by \eqref{eq:m-alloc}, so by monotonicity
(Observation~\ref{obs:m-gap}) and \eqref{eq:m-pef},
\begin{equation}
\label{eq:m-mono}
  \cost_i(A_j) \;\ge\; \cost_i(X_j) \;\ge\; \cost_i(X_i) .
\end{equation}
Note also that $T \subseteq P$, so $i \notin P$ implies $i \notin T$ and hence
$A_i = X_i$; likewise $j \notin P$ implies $A_j = X_j$ and $\subsidy_j = 0$.
Since $\agents$ is the disjoint union of $T$, $P \setminus T$ and
$\agents \setminus P$, the following three cases are exhaustive.

\smallskip
\noindent\textbf{Case 1: $i \in T$.} Then $\subsidy_i = 1$ and, by
Lemma~\ref{lem:m-terminal}(b), $\cost_i(A_i) = \cost_i(X_i) + \cost_i(e \mid X_i)
= \cost_i(X_i) + 1$, so $\cost_i(A_i) - \subsidy_i = \cost_i(X_i)$. Also
$i \in T \subseteq S$.
\begin{itemize}
  \item[(1a)] $j \in T$. Then $\subsidy_j = 1$ and $A_j = X_j \cup \set{e}$ for
    some $e \in R$, so $\cost_i(A_j) \ge \cost_i(X_i) + 1$ by
    Lemma~\ref{lem:m-expensive}. Hence
    $\cost_i(A_j) - \subsidy_j \ge \cost_i(X_i)$.
  \item[(1b)] $j \in P \setminus T$. Then $\subsidy_j = 1$ and $A_j = X_j$. By
    Lemma~\ref{lem:m-Pstructure}, $j \notin S$; since $i \in S$,
    Lemma~\ref{lem:m-terminal}(d) gives $\cost_i(X_i) < \cost_i(X_j)$, hence
    $\cost_i(X_j) \ge \cost_i(X_i) + 1$ by Observation~\ref{obs:m-gap}. Hence
    $\cost_i(A_j) - \subsidy_j = \cost_i(X_j) - 1 \ge \cost_i(X_i)$.
  \item[(1c)] $j \in \agents \setminus P$. Then $\subsidy_j = 0$ and
    $A_j = X_j$, so $\cost_i(A_j) - \subsidy_j = \cost_i(X_j) \ge \cost_i(X_i)$
    by \eqref{eq:m-pef}.
\end{itemize}

\smallskip
\noindent\textbf{Case 2: $i \in P \setminus T$.} Then $\subsidy_i = 1$ and
$A_i = X_i$, so $\cost_i(A_i) - \subsidy_i = \cost_i(X_i) - 1$. For every $j$
we have $\subsidy_j \le 1$ and $\cost_i(A_j) \ge \cost_i(X_i)$ by
\eqref{eq:m-mono}, hence
$\cost_i(A_j) - \subsidy_j \ge \cost_i(X_i) - 1 = \cost_i(A_i) - \subsidy_i$.

\smallskip
\noindent\textbf{Case 3: $i \in \agents \setminus P$.} Then $\subsidy_i = 0$
and $A_i = X_i$, so $\cost_i(A_i) - \subsidy_i = \cost_i(X_i)$.
\begin{itemize}
  \item[(3a)] $j \in \agents \setminus P$. Then $\subsidy_j = 0$ and
    $A_j = X_j$, so $\cost_i(A_j) - \subsidy_j = \cost_i(X_j) \ge \cost_i(X_i)$
    by \eqref{eq:m-pef}.
  \item[(3b)] $j \in T$. Then $\subsidy_j = 1$ and $A_j = X_j \cup \set e$ for
    some $e \in R$; we claim $\cost_i(A_j) \ge \cost_i(X_i) + 1$. If $i \in S$
    this is Lemma~\ref{lem:m-expensive}. If $i \notin S$, suppose
    $\cost_i(X_i) = \cost_i(X_j)$. Then $i \ne j$, since $j \in T \subseteq P$
    while $i \notin P$; so $(i,j) \in E$ by \eqref{eq:m-graph}, and $j \in P$,
    so (P2) of Definition~\ref{def:m-P} places $i \in P$ --- contradicting
    $i \notin P$. Hence $\cost_i(X_i) < \cost_i(X_j)$, so
    $\cost_i(X_j) \ge \cost_i(X_i) + 1$ by Observation~\ref{obs:m-gap}, and
    $\cost_i(A_j) \ge \cost_i(X_j)$ by monotonicity. In either case
    $\cost_i(A_j) - \subsidy_j \ge \cost_i(X_i)$.
  \item[(3c)] $j \in P \setminus T$. Then $\subsidy_j = 1$ and $A_j = X_j$; we
    claim $\cost_i(X_j) \ge \cost_i(X_i) + 1$. If $i \notin S$, suppose
    $\cost_i(X_i) = \cost_i(X_j)$. Then $i \ne j$, since $j \in P$ while
    $i \notin P$; so $(i,j) \in E$, and (P2) places $i \in P$, a contradiction.
    Hence $\cost_i(X_i) < \cost_i(X_j)$. If instead $i \in S$, then $j \notin S$
    by Lemma~\ref{lem:m-Pstructure}, and Lemma~\ref{lem:m-terminal}(d) gives
    $\cost_i(X_i) < \cost_i(X_j)$. In either case
    Observation~\ref{obs:m-gap} yields $\cost_i(X_j) \ge \cost_i(X_i) + 1$, so
    $\cost_i(A_j) - \subsidy_j = \cost_i(X_j) - 1 \ge \cost_i(X_i)$.
\end{itemize}

\smallskip
In every sub-case $\cost_i(A_j) - \subsidy_j \ge \cost_i(X_i) = \cost_i(A_i) -
\subsidy_i$ in Cases 1 and 3, and the inequality was shown directly in Case 2.
\end{proof}

% --------------------------------------------------------------------------
\subsection{The main theorem}
\label{sec:m-main}

\begin{theorem}
\label{thm:m-main}
For every $n$ and every instance $\langle \agents, \items,
\set{v_i}_{i \in \agents} \rangle$ with negative dichotomous valuations there
exist a complete allocation $\alloc$ and a subsidy vector
$\subsidy \in \set{0,1}^n$ with $\sum_{i \in \agents} \subsidy_i \le n-1$ such
that $(\alloc,\subsidy)$ is an envy-free solution. Both are computable in
polynomial time given value-oracle access.
\end{theorem}

\begin{proof}
Pass to cost form, $\cost_i = -v_i$, which is dichotomous in the sense of
Definition~\ref{def:dichcost}. Run the algorithm of Theorem~\ref{thm:m-twyz}
and let $X$, $R$, $r$ and $S$ be as fixed in \Cref{sec:m-terminal}.

If $r = 0$ then $X$ is a complete allocation satisfying \eqref{eq:m-pef}, so
$(X, \mathbf 0)$ is an envy-free solution with $\sum_i \subsidy_i = 0 \le n-1$.

If $r \ge 1$, construct $T$, $\alloc$ and $\subsidy$ as in
\Cref{sec:m-completion}. Then $\alloc$ is a complete allocation,
Theorem~\ref{thm:m-completion} shows $(\alloc,\subsidy)$ is an envy-free
solution, and Lemma~\ref{lem:m-budget} gives $\subsidy \in \set{0,1}^n$ with
$\sum_i \subsidy_i \le n-1$.

For the running time, the algorithm of Theorem~\ref{thm:m-twyz} is polynomial
and returns $X$, $R$ and $S$. Constructing $\alloc$ from \eqref{eq:m-alloc}
takes $r \le n-1$ assignments. The graph $G$ of \eqref{eq:m-graph} has $n$
vertices and is determined by the $n^2$ values $\cost_i(X_j)$, one oracle call
each. Computing $P$ is a backward reachability search in $G$ from $T$
restricted to vertices outside $S$, which takes $O(n^2)$ time. Hence the work
after Theorem~\ref{thm:m-twyz} is polynomial, and so is the whole procedure.
\end{proof}

By Example~\ref{ex:lowerbound} the bound of Theorem~\ref{thm:m-main} is tight:
there are instances on which every complete allocation requires total subsidy
at least $n-1$.

\begin{remark}
\label{rem:m-scope}
Theorem~\ref{thm:m-main} uses Theorem~\ref{thm:m-twyz} together with the four
properties of its halting state recorded in Lemma~\ref{lem:m-terminal}, which
are read off the rules (R1)--(R3) of~\cite{TWYZ23} restated in
\Cref{sec:m-terminal}; the statement of Theorem~\ref{thm:m-twyz} alone does not
supply them. Envy-freeness is verified directly from \eqref{eq:m-ef} in
Theorem~\ref{thm:m-completion}, so the argument does not invoke
Theorem~\ref{thm:hs-characterisation} or Theorem~\ref{thm:hs-minsubsidy}, and
the subsidy exhibited need not be the pointwise-minimal one. Part (a) of
Lemma~\ref{lem:m-terminal} is used only in Lemma~\ref{lem:m-budget}: the
conclusion of Theorem~\ref{thm:m-completion} holds whenever $r \le \abs S$.
\end{remark}
